
\section{第~8~章的练习}
\label{sec:ch14-exercises-ch8}

\vkexsync{44}%

\begin{exercise}\label{ex:ch14-044-gamma-kernel-infinite}
证明：\cref{thm:transversality-property}中所考虑的算子 $\Gamma$ 的\trKernel{}不是有限维的。
\end{exercise}

\begin{exercise}[定理~8.6.11~的另一证明]\label{ex:ch14-045-fredholm-another-proof}
设 $\Sigma\colon\RR\longrightarrow\operatorname{End}(\RR^{2n})$ 是一个连续映射，满足
\[
\Sigma(s)=\pi\,\mathrm{Id}\ \text{（当 } s<-s_0\text{）}
\qquad\text{且}\qquad
\Sigma(s)=3\pi\,\mathrm{Id}\ \text{（当 } s>s_0\text{）。}
\]
证明：\footnote{译注：原书此处第三项印作 $\sigma\cdot Y$，而 $\sigma$ 并未定义，上文刚定义的映射记作 $\Sigma(s)$；疑为大小写排版之误，应为 $\Sigma(s)\cdot Y$；此处径改。——译者注}由
\[
F(Y)=\frac{\partial Y}{\partial s}+J_0\,\frac{\partial Y}{\partial t}+\Sigma(s)\cdot Y
\]
所定义的算子 $F$ 是一个从 $W^{1,p}(\RR\times S^{1};\RR^{2n})$ 到 $L^{p}(\RR\times S^{1};\RR^{2n})$ 的\trFredholmOperator。
证明：
\[
\dim\operatorname{Ker}F=2n\,\#\,\bigl\{\ell\in\ZZ^{\star}\ \bigm|\ 1<2\ell<3\bigr\}=2n\text{，}
\]
然后再证明 $\dim\operatorname{Ker}F^{\star}=0$ 以及 $F$ 是\trSurjective{}的。
\end{exercise}
